\documentclass[11pt,a4paper]{article}
\usepackage[utf8]{inputenc}
\usepackage[vietnamese]{babel}
\usepackage{amsmath,amssymb}
\usepackage{geometry}
\geometry{a4paper, total={180mm,260mm}, left=15mm, right=15mm, top=18mm, bottom=18mm}

\usepackage{tikz}
\usetikzlibrary{shapes.geometric, arrows.meta, positioning, calc}
\tikzset{
  arr/.style={-{Stealth[scale=0.8]}},
  every node/.append style={align=center}
}
\usepackage{tcolorbox}
\tcbuselibrary{skins,breakable}
\usepackage{fontawesome}
\usepackage{fancyhdr}
\usepackage{enumitem}

\pagestyle{fancy}
\fancyhf{}
\renewcommand{\headrulewidth}{0.5pt}
\renewcommand{\footrulewidth}{0.5pt}
\fancyhead[L]{\small\bfseries\color{violet!85!black}\faBook\ SỔ TAY KIẾN THỨC TOÁN 2}
\fancyhead[R]{\small\bfseries\color{violet!85!black}Lớp Toán Thầy Lĩnh -- loptoanthaylinh.vn}
\fancyfoot[C]{\small\color{gray}-- Sổ tay kiến thức: Bài 48. Đơn vị, chục, trăm, nghìn --}

\setlength{\parindent}{0pt}
\setlength{\parskip}{0pt}

\begin{document}

% --- BANNER BÀI HỌC ---
\begin{tcolorbox}[
  enhanced,
  colback=violet!8!white,
  colframe=violet!70!black,
  arc=3mm, boxrule=1.2pt,
  left=3mm, right=3mm, top=2mm, bottom=2mm,
  halign=center
]
  {\LARGE\bfseries\color{violet!90!black}\faBookmark\ BÀI 48: ĐƠN VỊ, CHỤC, TRĂM, NGHÌN}\\[1mm]
  {\itshape\footnotesize\color{black!75} Quan hệ thập phân giữa các hàng: 10 đơn vị = 1 chục, 10 chục = 1 trăm, 10 trăm = 1 nghìn trên 1 trang A4}
\end{tcolorbox}

\vspace{-2mm}

\noindent
\begin{minipage}[t]{0.49\textwidth}\vspace{0pt}
\begin{tcolorbox}[
  enhanced,
  colback=violet!4!white,
  colframe=violet!75!black,
  arc=2mm, boxrule=0.8pt,
  left=2mm, right=2mm, top=1.8mm, bottom=1.8mm,
  title={\footnotesize\bfseries\color{white}\faKey\ 1. Mối quan hệ giữa các hàng}
]
{\footnotesize
\begin{itemize}\setlength{\itemsep}{1pt}\setlength{\leftskip}{-2.5mm}
  \item \textbf{Hệ đếm thập phân:}\\
  Cứ $10$ đơn vị ở hàng dưới thì hợp thành $1$ đơn vị ở hàng liền trên kề trước nó:
  \begin{itemize}\setlength{\itemsep}{0.5pt}\setlength{\leftskip}{-2.5mm}
    \item $10\text{ đơn vị} = \mathbf{1\text{ chục}}$ ($10$).
    \item $10\text{ chục} = \mathbf{1\text{ trăm}}$ ($100$).
    \item $10\text{ trăm} = \mathbf{1\text{ nghìn}}$ ($1\,000$).
  \end{itemize}
  \item \textbf{Số một nghìn:} Viết là $\mathbf{1\,000}$ (gồm chữ số $1$ và $3$ chữ số $0$ đứng sau; là số có $4$ chữ số).
\end{itemize}
}
\end{tcolorbox}
\end{minipage}
\hfill
\begin{minipage}[t]{0.49\textwidth}\vspace{0pt}
\begin{tcolorbox}[
  enhanced,
  colback=teal!4!white,
  colframe=teal!75!black,
  arc=2mm, boxrule=0.8pt,
  left=2mm, right=2mm, top=1.8mm, bottom=1.8mm,
  title={\footnotesize\bfseries\color{white}\faSliders\ 2. Kĩ năng đếm theo nhóm}
]
{\footnotesize
\begin{itemize}\setlength{\itemsep}{1pt}\setlength{\leftskip}{-2.5mm}
  \item \textbf{Đếm theo chục:} $10, 20, 30, 40, \dots, 100$ ($10$ chục là một trăm).
  \item \textbf{Đếm theo trăm:} $100, 200, 300, 400, 500, 600, 700, 800, 900, 1\,000$ ($10$ trăm là một nghìn).
  \item \textbf{Mô hình khối lập phương:}\\
  - $1$ ô vuông con $= 1$ đơn vị.\\
  - $1$ thanh gồm $10$ ô con $= 1$ chục.\\
  - $1$ tấm bảng gồm $10$ thanh $= 1$ trăm.\\
  - $1$ khối lớn gồm $10$ tấm bảng $= 1$ nghìn.
\end{itemize}
}
\end{tcolorbox}
\end{minipage}

\vspace{-1.5mm}
\begin{tcolorbox}[
  enhanced,
  colback=orange!4!white,
  colframe=orange!80!black,
  arc=2mm, boxrule=0.8pt,
  left=2mm, right=2mm, top=1.5mm, bottom=1.5mm,
  title={\footnotesize\bfseries\color{white}\faSitemap\ 3. Sơ đồ tư duy: Bậc thang chuyển đổi hệ đếm thập phân}
]
\centering
\begin{tikzpicture}[x=1.8cm, y=0.8cm]
  \node[rectangle, rounded corners=2pt, draw=violet!80!black, fill=violet!15, font=\scriptsize\bfseries] (dv) at (0,0) {ĐƠN VỊ\\($1$)};
  \node[rectangle, rounded corners=2pt, draw=teal!80!black, fill=teal!15, font=\scriptsize\bfseries] (ch) at (1.8,0) {CHỤC\\($10$)};
  \node[rectangle, rounded corners=2pt, draw=orange!80!black, fill=orange!15, font=\scriptsize\bfseries] (tr) at (3.6,0) {TRĂM\\($100$)};
  \node[rectangle, rounded corners=2pt, draw=red!80!black, fill=red!15, font=\scriptsize\bfseries] (ng) at (5.4,0) {NGHÌN\\($1\,000$)};

  \draw[->, thick, violet!80!black] (dv) -- node[above, font=\tiny\bfseries] {$\times 10$} (ch);
  \draw[->, thick, teal!80!black] (ch) -- node[above, font=\tiny\bfseries] {$\times 10$} (tr);
  \draw[->, thick, orange!80!black] (tr) -- node[above, font=\tiny\bfseries] {$\times 10$} (ng);
\end{tikzpicture}
\end{tcolorbox}

\vspace{-1.5mm}
\noindent
\begin{minipage}[t]{0.49\textwidth}\vspace{0pt}
\begin{tcolorbox}[
  enhanced,
  colback=green!4!white,
  colframe=green!60!black,
  arc=2mm, boxrule=0.8pt,
  left=2mm, right=2mm, top=1.8mm, bottom=1.8mm,
  title={\footnotesize\bfseries\color{white}\faCheckCircle\ 4. Ví dụ thực hành chuẩn mực}
]
{\footnotesize
\textbf{Điền số thích hợp vào chỗ chấm:}\\
- $30\text{ chục} = \mathbf{300}$ (ba trăm).\\
- $70\text{ đơn vị} = \mathbf{7}\text{ chục}$.\\
- $10\text{ trăm} = \mathbf{1\,000}$ (một nghìn).\\
- $5\text{ trăm} = \mathbf{50}\text{ chục} = \mathbf{500}\text{ đơn vị}$.
}
\end{tcolorbox}
\end{minipage}
\hfill
\begin{minipage}[t]{0.49\textwidth}\vspace{0pt}
\begin{tcolorbox}[
  enhanced,
  colback=red!4!white,
  colframe=red!75!black,
  arc=2mm, boxrule=0.8pt,
  left=2mm, right=2mm, top=1.8mm, bottom=1.8mm,
  title={\footnotesize\bfseries\color{white}\faExclamationTriangle\ 5. Mẹo tránh bẫy trắc nghiệm}
]
{\footnotesize
\begin{itemize}\setlength{\itemsep}{1pt}\setlength{\leftskip}{-2.5mm}
  \item \textbf{Bẫy đếm chữ số 0 của số 1000:}\\
  Số $1\,000$ có đúng \textbf{ba chữ số 0} đứng sau chữ số $1$. Nếu viết $100$ là một trăm, viết $10\,000$ là mười nghìn.
  \item Số bé nhất có 4 chữ số là $1\,000$.
\end{itemize}
}
\end{tcolorbox}
\end{minipage}


\end{document}
